\label{sec:results}

\subsection{Dataset Statistics}

The synthetic dataset comprises 6{,}000 sessions (5{,}000 benign, 1{,}000 malicious)
and 89{,}288 individual events. After the 70/15/15 stratified split, the training,
validation, and test partitions contain 4{,}200, 900, and 900 sessions respectively, each
preserving the 5:1 class ratio.

\subsection{Feature Extraction Results}

The feature extraction pipeline produces 24-dimensional session vectors (14 event-count
features, 4 time-delta statistics, 3 cardinality features, and 1 entropy feature; rare-event
flag columns are generated dynamically and are absent when all events appear in $\geq$5\% of
benign sessions, as in this corpus). The most discriminative features, as measured by Random
Forest Gini importance, are shown in Table~\ref{tab:features}.

\begin{table}[!t]
  \caption{Top-5 Feature Importances (mean Gini decrease)}
  \label{tab:features}
  \centering
  \begin{tabular}{lcc}
    \hline
    Feature & Importance & Rank \\
    \hline
    \texttt{td\_max} (max inter-event gap)  & 0.180 & 1 \\
    \texttt{td\_std} (gap std deviation)    & 0.135 & 2 \\
    \texttt{td\_min} (min inter-event gap)  & 0.120 & 3 \\
    \texttt{count\_4663} (object accesses)  & 0.112 & 4 \\
    \texttt{td\_mean} (mean inter-event gap)& 0.109 & 5 \\
    \hline
  \end{tabular}
\end{table}

\subsection{Classification Performance}

Table~\ref{tab:results} compares the proposed Random Forest against five baselines on the
held-out test set. All supervised baselines share the same 24-dimensional feature set and
imbalance handling; Isolation Forest is unsupervised (no labels used during training).

\begin{table}[!t]
  \caption{Classification Performance on Held-Out Test Set}
  \label{tab:results}
  \centering
  \begin{tabular}{lcccc}
    \hline
    Method & Precision & Recall & F1 (macro) & ROC-AUC \\
    \hline
    IF (unsupervised)      & 0.611 & 0.594 & 0.601 & 0.735 \\
    LR                     & 0.917 & 0.974 & 0.942 & 0.994 \\
    SVM (RBF)              & 0.818 & 0.881 & 0.844 & 0.951 \\
    Decision Tree          & 0.982 & 0.987 & 0.984 & 0.988 \\
    XGBoost                & 0.997 & 0.999 & 0.998 & 1.000 \\
    \textbf{RF (proposed)} & \textbf{1.000} & \textbf{1.000} & \textbf{0.999} & \textbf{1.000} \\
    \hline
  \end{tabular}
\end{table}

Random Forest achieves F1~=~0.999 (mean across 5 seeds, $\sigma$~=~0.001), the
highest of all supervised models. The improvement over Decision Tree (0.984) demonstrates
the variance-reduction benefit of ensembling. XGBoost is competitive (0.998) but requires
manual class-weight calibration; Random Forest handles imbalance via \texttt{class\_weight='balanced'}
with no additional tuning. The unsupervised Isolation Forest (0.601) establishes a
no-labels lower bound relevant to zero-day deployment scenarios. Time-delta statistics are
the dominant feature group; removing them degrades F1 by $-0.028$, while removing event
counts degrades by $-0.056$.
